\section{实验过程记录}

\subsection{ALU模块设计与仿真}

\textbf{工作内容：}
\begin{enumerate}
    \item 根据实验要求，使用Verilog语言编写32位ALU模块，支持加法、减法、与、或、非、SLT共6种运算
    \item 采用always@(*)块配合case语句实现组合逻辑，SLT运算使用\$signed()进行有符号比较
    \item 编写ALU测试激励，将B端口固定为32'h01，A端口依次输入不同的测试值
    \item 在Vivado中进行行为仿真，验证6种运算功能的正确性
    \item 整理仿真波形，记录验证结果
\end{enumerate}

\subsection{流水线全加器设计与仿真}

\textbf{工作内容：}
\begin{enumerate}
    \item 设计4级流水线32位全加器，将32位加法划分为4个流水阶段，每级处理8位
    \item 采用移位寄存器方案对齐各级部分和输出：第1级结果延迟3拍、第2级延迟2拍、第3级延迟1拍，最后与第4级结果拼接为32位输出
    \item 实现stall和flush控制信号。stall有效时所有寄存器保持；flush有效时清空全部寄存器，且flush优先级高于stall
    \item 编写测试激励，前10个周期进行多组加法运算（0x11110000+0x22220000至0x05550000+0x06660000），第10周期后stall 2个周期，第15周期flush 1个周期，之后验证0x64+0xC8=0x12C
    \item 在Vivado中进行行为仿真，整理仿真波形
\end{enumerate}

\subsection{问题与解决方案}

\subsection{问题1：SLT运算的符号位处理}
\textbf{问题描述：}
SLT指令需要将有符号数进行比较，但ALU输入为32位无符号类型。直接比较会将所有输入视为无符号数，导致比较结果错误。

\textbf{解决方案：}
使用Verilog内置函数\texttt{\$signed()}将操作数A和B转换为有符号数后再进行比较，确保SLT按照补码规则正确比较大小。当A < B时输出1，否则输出0。

\subsection{问题2：流水线暂停时数据保持}
\textbf{问题描述：}
当流水线暂停信号stall有效时，需要保持所有流水线级寄存器的值不变。如果在stall期间仍有新数据到达输入端，可能导致数据丢失或覆盖。

\textbf{解决方案：}
在所有流水线寄存器的always块中加入stall判断条件。当stall有效时，寄存器保持原值不更新；只有stall无效时才正常推进数据。

\subsection{问题3：flush与stall的优先级}
\textbf{问题描述：}
当flush和stall同时有效时，需要确定哪个信号优先级更高。

\textbf{解决方案：}
flush信号用于处理分支预测失败等异常情况，需要立即清空流水线，因此flush优先级高于stall。在Verilog代码中，优先判断flush信号，只有在flush无效时才判断stall。

\subsection{问题4：流水线各级部分和的对齐}
\textbf{问题描述：}
4级流水线中，各级的部分和在不同时钟周期产生。如果直接将各级结果拼接，会出现不同批次数据的部分和混在一起，导致输出错误。

\textbf{解决方案：}
采用移位寄存器方案。将第1级部分和依次经过3级移位寄存器，第2级经过2级，第3级经过1级，使所有部分和在同一时钟周期到达输出端。同时，各阶段的剩余操作数也需逐级传递，确保每级处理正确的数据段。
